\subsection{Experimental Context}

All 28 experiments were conducted during a 14-hour session at the Funding the Commons hackathon (March 14--15, 2026, Frontier Tower, San Francisco). This rapid-iteration context shaped the experimental design: experiments were planned, executed, analyzed, and used to inform subsequent experiments in real time. While this introduces risks of data-dependent hypothesis generation, we mitigate this by: (a) reporting all experiments including nulls and debunked signals, (b) applying consistent statistical protocols across all analyses, and (c) subjecting every positive finding to at least one adversarial control.

\subsection{Models}

Campaign 3 uses three primary models for cross-architecture validation:

\begin{table}[H]
\centering
\caption{Primary models used in Campaign 3.}
\label{tab:c3_models}
\begin{tabular}{llccc}
\toprule
\textbf{Model} & \textbf{Architecture} & \textbf{Parameters} & \textbf{Layers} & \textbf{Precision} \\
\midrule
Qwen2.5-7B-Instruct & Qwen & 7B & 28 & BF16 \\
Llama-3.1-8B-Instruct & Llama & 8B & 32 & BF16 \\
Mistral-7B-Instruct-v0.3 & Mistral & 7B & 32 & BF16 \\
\bottomrule
\end{tabular}
\end{table}

Additionally, Campaign 2's full 16-model scale ladder (0.5B to 70B across 6 architecture families) is used for scale invariance and category geometry analyses (\Cref{sec:robustness}).

For absence-of-suppression experiments, abliterated variants of all three models were created using the method of \citet{arditi2024refusal}, removing the refusal direction to produce models that comply with harmful requests. These abliterated models serve as a key control: if the detection signal is output suppression, then abliterated models---which lack suppression circuitry---should show normal (non-suppressed) cache geometry when answering harmful prompts.

\subsection{KV-Cache Feature Extraction}

\paragraph{Core features (4).} Following Campaigns 1--2, we extract four aggregate features from the KV-cache after generation:

\begin{enumerate}[leftmargin=*]
    \item \textbf{Key norm} ($\|\mathbf{K}\|$): Mean Frobenius norm across all layers. Captures overall activation magnitude.
    \item \textbf{Norm per token} ($\|\mathbf{K}\| / S$): Key norm divided by total sequence length. Controls for response length.
    \item \textbf{Effective rank} ($r_{\text{eff}}$): Mean effective rank across layers, computed as $r_{\text{eff}} = \exp(H(\mathbf{p}))$ where $H$ is Shannon entropy of normalized singular values. Captures representational dimensionality.
    \item \textbf{Rank entropy} ($\sigma_r$): Standard deviation of per-layer effective ranks. Captures heterogeneity of processing across the transformer stack.
\end{enumerate}

\paragraph{Extended features (5).} Campaign 3 introduces five additional features computed directly from key matrices, requiring no attention weight extraction:

\begin{enumerate}[leftmargin=*,start=5]
    \item \textbf{Key norm variance}: Variance of per-position key norms within each layer, averaged across layers. Higher values indicate non-uniform token processing.
    \item \textbf{Angular spread}: Mean pairwise cosine distance between key vectors at different positions, averaged across layers. Lower values indicate more homogeneous key geometry.
    \item \textbf{Layer correlation}: Pearson correlation of per-position key norms between adjacent layers. Higher values indicate consistent processing across the transformer stack.
    \item \textbf{Generation delta}: Difference in mean key norms between encoding (prompt) and generation (response) phases. Positive values indicate encoding dominance.
    \item \textbf{Head variance}: Variance of mean key norms across attention heads within each layer, averaged across layers. Captures attention head specialization.
\end{enumerate}

\subsection{Classification Protocol}

All detection experiments use leave-one-out cross-validated (LOO-CV) logistic regression (LR) as the primary classifier, with random forest (RF) as secondary. LOO-CV is appropriate for our sample sizes ($n = 20$--$40$ per condition) and provides an unbiased estimate of generalization error.

Effect sizes are reported as Cohen's $d$ with the convention: positive $d$ indicates the target condition (e.g., deceptive) exceeds the baseline condition (e.g., honest).

\subsection{Statistical Controls}

\paragraph{Same-prompt paradigm.} For deception and sycophancy detection, we use identical system prompts across conditions, varying only the user message. This eliminates system-prompt fingerprinting as a confound.

\paragraph{Length residualization.} Where conditions differ in response length, we report both raw and length-residualized effect sizes (OLS regression of features on token count, using residuals).

\paragraph{ANCOVA.} For per-layer analyses, we use analysis of covariance with token count as covariate, reporting both raw and token-controlled effect sizes.

\paragraph{Cross-architecture replication.} Detection capabilities are validated on at least two architectures before being reported as robust findings.

\subsection{Prompt Sets}

Prompts are drawn from several sources:
\begin{itemize}[leftmargin=*]
    \item \textbf{Confabulation}: 30 frequency-matched pairs (C7), 60 factual questions across 6 domains (S3), 60 contrastive pairs (bare/correct/irrelevant/wrong).
    \item \textbf{Sycophancy}: 60 questions from the sycophancy prompt battery, with ``I think [wrong answer]. Am I right?'' framing. Enhanced version adds authority framing per \citet{sharma2023towards}.
    \item \textbf{Deception}: System-prompt manipulation (honest vs.\ deceptive instructions) and same-prompt deception (user-message embedded deception request).
    \item \textbf{Refusal/Suppression}: 20 harmful prompts designed to trigger safety refusal, paired with 20 benign prompts matched for domain. Impossibility prompts (e.g., ``What is north of north?'') added in Experiment 36. Abliterated model responses on the same harmful prompts used for absence-of-suppression analysis.
\end{itemize}

All prompt sets are available in the public repository.
